%==============================================================
\section{Online Newton estimation: regret for $K=1$}\label{sec:regret}
%==============================================================

\subsection{Algorithm and the basic regret reduction}

The estimator is event-indexed. With $F_0\succ0$ a prior weight,
\[
\Delta F_i=\int_0^{w_i}\frac{\phi\phi\T}{\lambda_{\theta_i}}\dd u,\quad F_i=F_{i-1}+\Delta F_i,\quad
g_i=b_i-\frac{a_i}{a_i\T\theta_i},\quad \theta_{i+1}=\Pi^{F_i}_{\Theta}\big(\theta_i-F_i^{-1}g_i\big),
\]
where $\Pi^{F}_\Theta$ is the projection in the norm $\norm x_F^2=x\T Fx$ and $\Theta=\Theta_c$ is a compact box containing $\theta^*$ (Proposition~\ref{prop:projbox}).
Put $v_i=\theta_i-\theta^*$, $P_0=\norm{v_1}_{F_0}^2$ and the \emph{pseudo-regret}
$R_n=\sum_i\int_0^{w_i}\lambda^*\psi(\lambda_{\theta_i}/\lambda^*)\dd u$, $\psi(r)=r-1-\log r$ of Definition~\ref{def:regret}; by Lemma~\ref{lem:regid}, $\E R_n=\mathrm{Reg}_n$, the regret of Section~\ref{sec:regretnotion}.

\begin{proposition}[regret reduction]\label{prop:reg}\Pt{}
Let $\theta_i$ be $\cF_{t_{i-1}}$-measurable, $\Theta$ convex, $\theta^*\in\Theta$. Then
\begin{enumerate}[label=(\alph*),nosep]
\item $\E[\ip{g_i}{v_i}\mid\text{past}]=\E[v_i\T\Delta F_iv_i\mid\text{past}]$;
\item $\E[\text{increment of }R\mid\text{past}]\le\E[\ip{g_i}{v_i}\mid\text{past}]$ (since $\psi(r)\le(r-1)^2/r$);
\item $\E R_n\le X:=\E\sum_iv_i\T\Delta F_iv_i\le P_0+G$, $G:=\E\sum_ig_i\T F_i^{-1}g_i$.
\end{enumerate}
No exp-concavity is used.
\end{proposition}
\begin{proof}
(a) Given the past, $\phi,\lambda^*,\lambda_\theta$ on the gap are deterministic functions of $u$, with $\P(\tau>u)=e^{-\Lambda^*(u)}$. Since $\phi\T v=\lambda_\theta-\lambda^*$, $\ip gv=\int_0^\tau(\lambda_\theta-\lambda^*)\dd u-(\lambda_\theta-\lambda^*)(\tau)/\lambda_\theta(\tau)$; by Fubini its mean is $\int e^{-\Lambda^*}(\lambda_\theta-\lambda^*)^2/\lambda_\theta\dd u=\E\int_0^\tau(\phi\T v)^2/\lambda_\theta\dd u$.
(b) $\log r\ge1-1/r$. (c) By non-expansiveness of $\Pi^{F_i}$ and $\theta^*\in\Theta$: $\norm{v_{i+1}}_{F_i}^2\le\norm{v_i-F_i^{-1}g_i}_{F_i}^2$, i.e.\ $2\ip{g_i}{v_i}\le\norm{v_i}_{F_i}^2-\norm{v_{i+1}}_{F_i}^2+g_i\T F_i^{-1}g_i$. Sum, using $\norm{v_i}_{F_i}^2-\norm{v_i}_{F_{i-1}}^2=v_i\T\Delta F_iv_i$: $2\sum\ip gv\le P_0+\sum v\T\Delta Fv+\sum g\T F^{-1}g$. Take expectations and use (a): $2X\le P_0+X+G$.
\end{proof}
\Nt{}: Newton identity residual $\le7\times10^{-13}$; $\E\sum\ip gv=11.46\pm0.25$ vs $\E\sum v\T\Delta Fv=11.40\pm0.47$.

\begin{proposition}[ratio set, deterministic]\label{prop:ratio}\Pt{}
If $\Theta_c=[\mu_{lo},\mu_{hi}]\times[\alpha_{lo},\alpha_{hi}]$ with $\alpha_{hi}\le\beta(1-\eps/2)$, $c^+=\max(\mu_{hi}/\mu^*,\alpha_{hi}/\alpha^*)$, $c^-=\min(\mu_{lo}/\mu^*,\alpha_{lo}/\alpha^*)$, then $c^-\lambda^*\le\lambda_\theta\le c^+\lambda^*$ pathwise for every $\theta\in\Theta_c$ (the ratio $(\mu+\alpha s)/(\mu^*+\alpha^*s)$ lies between $\mu/\mu^*$ and $\alpha/\alpha^*$). Hence $F_{i-1}\succeq F_0+J_{t_{i-1}}/c^+$ where $J_T=\int_0^Taa\T/\lambda^*\dd t$ is the oracle information.
\end{proposition}
This removes the circularity (the ratio-set assumption holds by design), at the price that $c^\pm$ are box constants.

\subsection{Score--drift decomposition and the conditional bound}

Write $g_i=s_i+h_i$ with $s_i=b_i-a_i/\lambda^*(\tau_i)$ (score at $\theta^*$) and $h_i=a_i/\lambda^*(\tau_i)-a_i/\lambda_\theta(\tau_i)=a_i(a_i\T v_i)/(\lambda^*\lambda_\theta)(\tau_i)$.
Let $\Delta J_i=\int_0^{w_i}\phi\phi\T/\lambda^*\dd u$ and $\ell_i$ below denote the \emph{point leverage}
\[
\bar\ell_i\ \ge\ \sup_{y\in[0,S_{t_{i-1}}]}\ \frac{(1,y)F_{i-1}^{-1}(1,y)\T}{\mu^*+\alpha^*y}.
\]

\begin{lemma}[score is centred, covariance $=\Delta J$]\label{lem:score}\Pt{}
$\E[s_i\mid\text{past}]=0$, $\E[s_is_i\T\mid\text{past}]=\E[\Delta J_i\mid\text{past}]$.
(The stopped martingale $M_t=\int_0^t(\phi/\lambda^*)(\dd N-\lambda^*\dd u)$ has $M_\tau=-s_i$ and $\langle M\rangle_\tau=\Delta J_i$.)
\end{lemma}

\begin{lemma}[drift bound]\label{lem:drift}\Pt{}
If $\lambda_\theta\ge\mu_{lo}$ and $\bar\ell_i$ bounds the point leverage, then $\E[h_i\T F_{i-1}^{-1}h_i\mid\text{past}]\le(\bar\ell_i/\mu_{lo})\E[v_i\T\Delta F_iv_i\mid\text{past}]$.
\end{lemma}
\begin{proof}
$h\T F^{-1}h=[\phi\T F^{-1}\phi/\lambda^*](\tau)\cdot(\phi\T v)^2/(\lambda^*\lambda_\theta^2)(\tau)$; with density $\lambda^*e^{-\Lambda^*}$ of $\tau$ the mean is $\int e^{-\Lambda^*}[\phi\T F^{-1}\phi/\lambda^*](\phi\T v)^2/\lambda_\theta^2\dd u\le(\bar\ell/\mu_{lo})\E\int(\phi\T v)^2/\lambda_\theta\dd u$.
\end{proof}

\begin{lemma}[score part and $\log\det$, truncated]\label{lem:logdet}\Pt{}
Let $E_i=\Lambda^*(\tau_i)$ (conditionally $\mathrm{Exp}(1)$ and independent of the past by the time-change theorem), $q_i:=\tr(F_{i-1}^{-1}\E[\Delta J_i\mid\text{past}])=\E[s_i\T F_{i-1}^{-1}s_i\mid\text{past}]$, and assume $\lambda_\theta\le c^+\lambda^*$, $\bar\ell_i\le\bar\ell$. For $m>0$ put $\bar x=\bar\ell m/(c^-\mu^*)$, $r(x)=x/\log(1+x)$, $\ell'_i=\log\det(F_i/F_{i-1})\1\{E_i\le m\}$. Then
\[
q_i\ \le\ c^+r(\bar x)\,\E[\ell'_i\mid\text{past}]+\frac{\bar\ell}{\mu^*}(m+1)e^{-m}.
\]
\end{lemma}
\begin{proof}
On $\{E_i\le m\}$, $\tr(F^{-1}\Delta F)=\int\phi\T F^{-1}\phi/\lambda_\theta\dd u\le\bar\ell\tau/c^-\le\bar\ell E_i/(c^-\mu^*)\le\bar x$, so every eigenvalue $y_j$ of $F^{-1/2}\Delta FF^{-1/2}$ lies in $[0,\bar x]$ and $\sum y_j\le r(\bar x)\sum\log(1+y_j)$ by concavity; also $\Delta J\preceq c^+\Delta F$. On $\{E_i>m\}$, $\tr(F^{-1}\Delta J_i)\le\bar\ell E_i/\mu^*$ and $\E[E;E>m]=(m+1)e^{-m}$.
\end{proof}
\begin{remark}[correction to the earlier Theorem on conditional regret]
The earlier assumption ``gap leverage $x_i\le\bar x$ almost surely'' cannot hold because $\tr(F^{-1}\Delta F)\propto E_i$ is unbounded. The truncation above is the repaired form and costs nothing in the leading constant.
\end{remark}

\begin{theorem}[conditional regret bound]\label{thm:cond}\Ps{}
Assume a.s.\ for $i\le n$: $c^-\lambda^*\le\lambda_{\theta_i}\le c^+\lambda^*$, $\lambda_{\theta_i}\ge\mu_{lo}$, $\bar\ell_i\le\bar\ell$ and $B:=(1+1/\eta)\bar\ell/\mu_{lo}<1$. Then, for every $m>0$,
\[
\E R_n\ \le\ \frac{P_0+(1+\eta)\big[c^+r(\bar x)\,\E\logdet(F_n/F_0)+n\frac{\bar\ell}{\mu^*}(m+1)e^{-m}\big]}{1-B}.
\]
(Take $m=\ln n+x$ so that the last term is $e^{-x}$-small.) Proof: Proposition~\ref{prop:reg}(c), $g\T F^{-1}g\le(1+\eta)s\T F^{-1}s+(1+1/\eta)h\T F^{-1}h$, Lemma~\ref{lem:score}, Lemma~\ref{lem:logdet}, Lemma~\ref{lem:drift}, and solve $X\le P_0+AD+BX$.
\end{theorem}
For $K=1$, $\E\logdet(F_n/F_0)\approx2\log T+\logdet I-\logdet F_0$, so the bound has the right growth in $T$ with a constant $\gtrsim A\ge1.44$ and a further factor $2$ lost in $\psi\le(r-1)^2/r$.
\Nt{}: Lemma~\ref{lem:score}: $\E\sum s\T F^{-1}s=6.91$ vs $\E\sum\tr(F^{-1}\Delta J)=6.77$; Lemma~\ref{lem:drift}: $0.024\le3.4$ (the drift term is empirically negligible).
The point leverage requires about $4(k-1)/(\eps^2\mu_{lo})$ events before $2\bar\ell/\mu_{lo}<1$ (observed $1520$ at $\eps=0.1$, $10086$ at $\eps=0.03$ against $2000$, $22000$ predicted): the leverage burn-in is $\asymp1/\eps^2$ events.

\subsection{Concentration of the accumulated information}

Let $\gamma=\beta\eps$, $\mathrm{ES}=\mu/(\beta\eps)$, $\sigma_S=(\E S^2+3(\E S)^2)^{1/2}$. Take two disjoint intervals $I_j=[c_j-d,c_j+d]$ and trapezoids $h_j:\R_+\to[0,1]$ supported in $I_j$ with ramps of width $\rho$ and plateau $[c_j-d/2,c_j+d/2]$ ($\rho=d/2$), centres $c_1=0.5\,\mathrm{ES}$, $c_2=1.8\,\mathrm{ES}$, $d=0.15\,\mathrm{ES}$.

\begin{lemma}[variance of an occupation functional]\label{lem:F1}\Pt{}
For $h\in[0,1]$ Lipschitz with constant $\Lambda_h$, $m=\E_\pi h$: $\Var(\int_0^Th(S_t)\dd t)\le2T\Lambda_h\sqrt m\,\sigma_S/\gamma$.
\end{lemma}
\begin{proof}
$|P_th(s)-m|\le\Lambda_he^{-\gamma t}(s+\mathrm{ES})$ by Lemma~\ref{lem:ellcontr}; Cauchy--Schwarz with $\Var h\le m$ gives $|\Cov(t)|\le\Lambda_h\sqrt m\sigma_Se^{-\gamma t}$ and $\Var\le2T\int_0^\infty|\Cov|$.
\end{proof}

\begin{lemma}[two-interval lower bound]\label{lem:F2}\Pt{}
Let $O_j=\int_0^T\1\{S_t\in I_j\}\dd t$, $L=\mu+\alpha(c_2+d)$ and $c_2-c_1>2\sqrt2\,d$. For all $v\in\R^2$:
$v\T J_Tv\ge(\min_jO_j/L)\,v\T Qv$, $Q=\begin{psmallmatrix}1&(c_1+c_2)/2\\(c_1+c_2)/2&(c_1^2+c_2^2)/2-2d^2\end{psmallmatrix}\succ0$.
\end{lemma}
\begin{proof}
On $I_j$: $v\T a=(v_0+v_1c_j)+v_1(s-c_j)$, so $(v\T a)^2\ge\frac12(v_0+v_1c_j)^2-v_1^2d^2$; divide by $\lambda^*\le L$ and sum with weights $O_j$; $\det Q=\Delta^2/4-2d^2>0$.
\end{proof}

\begin{assumption}[M]\label{ass:M}
$m:=\min_j\E_\pi h_j>0$. Numerically $m\approx0.15$ ($I_1$), $0.05$ ($I_2$) at $\eps=0.3,0.1,0.03$ ($k_0=2$). A rigorous lower bound, valid for all $\eps\le0.3$ tested and $k_0\in[0.6,4]$, is Proposition~\ref{prop:Mexact}; for $K\ge2$ it is replaced by Theorem~\ref{thm:FK}.
\end{assumption}

\begin{theorem}[concentration of $J_T$, Chebyshev tails]\label{thm:Fconc}\Ps{} (modulo Assumption~\ref{ass:M})
For $q\in(0,1)$ and $T_0\ge16\sigma_S/(\rho\gamma m^{3/2}q)$, with probability $\ge1-2q$, simultaneously for all $T\ge T_0$: $J_T\succeq(mT/(4L))Q$. (Chebyshev with Lemma~\ref{lem:F1} and a union over $T_k=2^kT_0$; $J_T$ is nondecreasing.)
\end{theorem}
Rigorous thresholds are $\sim10^4\times$ the empirical ones, but scale correctly ($T_0\propto1/\gamma$). Under $\mathrm{(L1)}$ it implies $F_{i-1}\succeq(mt_{i-1}/(4c^+L))Q$, hence point leverage $\bar\ell_i=O(c^+/(\eps mt))$.

\subsection{Exponential tails}

\begin{lemma}[exact identities]\label{lem:E1}\Pt{}
With $M'=N-\Lambda$: (a) $\eps\Lambda_T=\mu T+nM'_T-n(S_T-S_0)$; (b) $\int_0^T(S_t-\mathrm{ES})\dd t=(S_0-S_T)/\gamma+(N_T-\Lambda_T)/\gamma$.
\end{lemma}

\begin{lemma}[tails of the integrated intensity]\label{lem:E2}\Pt{}
$\P(\Lambda_T\le(1-\eta)\mu T/\eps,\ |S_T-S_0|\le y)\le\exp\!\big(-(\eta\mu T/n-y)^2\eps/(2\mu T)\big)$ (supermartingale $\exp(\theta M'-\theta^2\Lambda/2)$, $\theta<0$); the upper tail is Bernstein. Rate $\asymp\mu\eps T=(\mu/\beta)\gamma T$.
\end{lemma}

\begin{theorem}[exponential-tail occupation bound]\label{thm:E3}\Ps{}
For $h\in[0,1]$, $\Lambda_h$-Lipschitz, $m=\E_\pi h$, $\hat u(s)=\int_0^\infty(P_th(s)-m)\dd t$ satisfies $|\hat u(s)-\hat u(s')|\le\Lambda_h|s-s'|/\gamma$ and Dynkin gives $\int_0^T(h(S_t)-m)\dd t=\hat u(S_0)-\hat u(S_T)+\mathcal M_T$ with jumps of $\mathcal M$ bounded by $b_0=\Lambda_h/\gamma$ and $\langle\mathcal M\rangle_T\le b_0^2\Lambda_T$. Then for $T\gamma\gtrsim\kappa/m$ ($\kappa:=\Lambda_h\mathrm{ES}$)
$\P(\int_0^Th\le Tm/2)\le3e^{-cT\gamma}$, $c=m^2\mu/(32\kappa^2(1+\eta)\beta)$.
\end{theorem}
Integrability: $|\hat u(s)|\le\Lambda_h(s+\mathrm{ES})/\gamma$, $\E\Lambda_T<\infty$, so $\mathcal M$ is a square-integrable martingale and Dynkin applies. The rate $c\approx3\times10^{-7}$ is $10^5\times$ below the observed decay; this affects constants only.

\subsection{Bootstrap: the iterates enter the ratio set $R(\delta)=\{|\mu/\mu^*-1|,|\alpha/\alpha^*-1|\le\delta\}$}

\begin{lemma}[martingale structure]\label{lem:C1}\Pt{}
For the projected iteration and every restart index $i_0$ with $\cG_{i_0-1}$-measurable data,
\[
D_j\le D'_0-X_j-M_j+G_j,\qquad D_j=\norm{v_{j+1}}_{F_j}^2,\ X_j=\sum_{i=i_0}^jv_i\T\Delta F_iv_i,\ G_j=\sum g_i\T F_i^{-1}g_i,
\]
$M_j=\sum m_i$, $m_i=-2\int_0^{\tau_i}(\phi\T v_i/\lambda_\theta)(\dd N-\lambda^*\dd u)$. $M$ is a martingale with integrand $|w|\le b_M:=2\kappa_1$, $\kappa_1=\max(1,1/c^--1)$, and $\langle M\rangle\le(4/c^-)X$.
\end{lemma}

\begin{lemma}[exponential bound for the Newton martingale]\label{lem:E4}\Pt{}
For $\rho\in(0,1)$ and $0<\theta\le\min(1/b_M,\rho c^-/(4(e-2)))$: with probability $\ge1-e^{-x}$, for all $t$, $-M_t\le\rho X_t+x/\theta$. ($e^{-\theta w}-1+\theta w\le(e-2)\theta^2w^2$, so $\exp(-\theta M-(e-2)\theta^2\langle M\rangle)$ is a supermartingale; Ville.)
\end{lemma}

\paragraph{The gradient-norm term $G_n$ with exponential tail (``E5'').}
Write $Q_i=s_i\T F_{i-1}^{-1}s_i$, $H_i=h_i\T F_{i-1}^{-1}h_i$.

\begin{lemma}[pathwise bound on one score term]\label{lem:E5a}\Pt{}
$Q_i\le(2\bar\ell_i/\mu^*)(1+E_i^2)$.
\end{lemma}
\begin{proof}
$Q\le2b\T F^{-1}b+2a\T F^{-1}a/\lambda^{*2}$. Since $\lambda^*\ge\mu^*$, $a\T F^{-1}a/\lambda^{*2}\le\bar\ell/\mu^*$. For $b=\int\phi$: $\abs b_{F^{-1}}\le\int\sqrt{\phi\T F^{-1}\phi}\le\sqrt{\bar\ell}\int\sqrt{\lambda^*}\le\sqrt{\bar\ell\tau E}$ and $\tau\le E/\mu^*$.
\end{proof}
(Checked: $0$ violations in $1.44\times10^6$ simulated gaps, max $Q/\text{bound}=0.26$.)

\begin{lemma}[score part, exponential tail]\label{lem:E5b}\Pt{}
Let $\bar\ell_i\le\bar\ell$ on $[i_0,n]$, $\bar B:=(2\bar\ell/\mu^*)(1+m^2)$, $0<\eps'\le e-2$. With probability $\ge1-ne^{-m}-e^{-x}$, for all $j\in[i_0,n]$:
$\sum_{i_0}^jQ_i\le(1+\eps')\sum_{i_0}^jq_i+(e-2)\bar Bx/\eps'$.
\end{lemma}
\begin{proof}
On $\{E_i\le m\ \forall i\}$, $Q_i=Q_i':=\min(Q_i,\bar B)$ by Lemma~\ref{lem:E5a}; $Q'_i\in[0,\bar B]$, $\E[Q_i'|\text{past}]\le q_i$, $\E[Q_i'^2|\text{past}]\le\bar Bq_i$. With $e^y\le1+y+(e-2)y^2$ ($y\le1$), $\exp(\theta\sum(Q'-\E Q')-(e-2)\theta^2\bar B\sum q)$ is a supermartingale for $\theta\le1/\bar B$; take $\theta=\eps'/((e-2)\bar B)$ and apply Ville.
\end{proof}

\begin{lemma}[conversion of conditional means]\label{lem:E5c}\Pt{}
For $\ell'_i\in[0,b']$ and $0<\eps''\le1$: with probability $\ge1-e^{-x}$, for all $j$, $\sum\E[\ell'_i|\text{past}]\le(1-\eps''/2)^{-1}(\sum\ell'_i+b'x/\eps'')$ (from $\E e^{-\theta\ell'}\le\exp(-\theta(1-\theta b'/2)\E\ell')$ with $\theta=\eps''/b'$).
\end{lemma}

\begin{lemma}[drift part]\label{lem:E5d}\Pt{}
$H_i\le(\bar\ell_i/\mu_{lo})(a_i\T v_i)^2/(\lambda^*\lambda_\theta)(\tau_i)$ and $(a\T v)^2/(\lambda^*\lambda_\theta)(\tau_i)=v_i\T\Delta F_iv_i+\nu_i$, where $\nu_i$ is the increment of $N^A=\int w_A(\dd N-\lambda^*\dd u)$, $w_A=(\phi\T v)^2/(\lambda_\theta\lambda^*)\le b_A:=\kappa_1^2c^+$, $\langle N^A\rangle\le b_AX$. Ville with $\theta=1/b_A$: with probability $\ge1-e^{-x}$, $N^A_j\le(e-2)X_j+b_Ax$, so
$\sum_{i\le j}H_i\le(\bar\ell/\mu_{lo})\big((e-1)X_j+b_Ax\big)$.
\end{lemma}

\begin{proposition}[E5: exponential tail for $G_n$]\label{prop:E5}\Pt{}
Assume for $i\in[i_0,n]$: $c^-\lambda^*\le\lambda_\theta\le c^+\lambda^*$, $\lambda_\theta\ge\mu_{lo}$, $\bar\ell_i\le\bar\ell$ and $B_X:=(1+1/\eta)(e-1)\bar\ell/\mu_{lo}\le1/4$. With $m=\ln n+x$, with probability $\ge1-4e^{-x}$, for all $j\in[i_0,n]$:
\begin{multline*}
G_j\le(1+\eta)(1+\eps')\frac{c^+r(\bar x)}{1-\eps''/2}\Big(\logdet\frac{F_j}{F_{i_0-1}}+\frac{b'x}{\eps''}\Big)\\
+(1+\eta)\Big[(1+\eps')\frac{\bar\ell}{\mu^*}(m+1)e^{-x}+(e-2)\frac{\bar Bx}{\eps'}\Big]+B_XX_j+(1+\tfrac1\eta)\frac{\bar\ell}{\mu_{lo}}b_Ax.
\end{multline*}
\end{proposition}
\begin{proof}
$G_j\le\sum g\T F_{i-1}^{-1}g$ ($F_i\succeq F_{i-1}$), $g\T F^{-1}g\le(1+\eta)Q+(1+1/\eta)H$; apply Lemmas~\ref{lem:E5b}, \ref{lem:logdet}, \ref{lem:E5c}, \ref{lem:E5d}. Four failure events (truncation, three Ville events).
\end{proof}
When $\bar\ell(\ln n+x)^2\to0$ and $c^+\to1$ all coefficients of $x$ vanish: $G_j\le(1+o(1))\logdet(F_j/F_{i_0-1})+o(1)\,x+B_XX_j$.
Realised $\sum Q/\sum q$: median $0.95$--$1.02$ ($5$--$95\%$: $0.81$--$1.23$).

\begin{theorem}[bootstrap with exponential tails]\label{thm:Cexp}\Ps{}
With $\rho=1/2$ and $B_X\le1/4$, Lemma~\ref{lem:E4} and Proposition~\ref{prop:E5} give, with probability $\ge1-5e^{-x}$,
\[
\sup_{i_0\le j\le n}D_j+\tfrac14X_n\ \le\ D'_0+x/\theta_{1/2}+\mathcal G_n(x),
\]
$\mathcal G_n(x)$ the $X$-free part of Proposition~\ref{prop:E5}. Since $D_j\ge\lambda_{\min}(F_j)|v_{j+1}|^2$ and $\lambda_{\min}(F_j)\ge\lambda_T:=mT\lambda_Q/(4c^+L)$ on the event of Theorem~\ref{thm:Fconc}/\ref{thm:E3}, the iterates lie in $R(\delta)$ for $t\ge T_1(\delta,q)=O\!\big((\ln(1/q)+d\ln T)/(\gamma\delta^2)\big)$ with probability $\ge1-q$.
\end{theorem}
The restart index must be predictable: for a deterministic time $T_0$ use the first gap starting after $T_0$ (event $\{t_{i-1}\ge T_0\}\in\cG_{i-1}$).

\subsection{Expected regret with the sharp constant}

Define the \emph{good-gap} indicator $\1_{\mathrm{good},i}=\1\{\theta_i\in R(\delta)\}\1\{\lambda_{\min}(F_{i-1})\ge\lambda_{t_{i-1}}\}\1\{\bar\ell_i\le\bar\ell_{t_{i-1}}\}\1\{S_{t_{i-1}}\le s_{t_{i-1}}\}$. It is $\cG_{i-1}$-measurable (predictable). Multiplying the increments in Lemmas~\ref{lem:E5b}--\ref{lem:E5d} by a predictable indicator in $[0,1]$ keeps all the supermartingales supermartingales, so no stopping times or restarts are needed: one global telescoping (Lemma~\ref{lem:C1} with $i_0=1$) suffices, and $M_n$ is a square-integrable martingale with $\E M_n=0$ (no optional stopping needed).

\begin{theorem}[expected regret, sharp constant, $K=1$, $k_0>1$]\label{thm:R}\Ps{}
Under the standing assumptions (box $\Theta_c$, exact $F$-projection, Assumption~\ref{ass:M}) and for every $\delta,\eta>0$ there is $C=C(\delta,\eta,\eps,\mu,\beta,\text{box},F_0)$, independent of $n$, with
\[
\E R_n\ \le\ \big(\tfrac12+O(\delta+\eta)\big)\,d\,\ln T_n+C,\qquad d=2,\qquad\text{so}\quad \limsup_n\frac{\E R_n}{\ln T_n}\le\frac d2=1 .
\]
\end{theorem}
\begin{proof}
Fix $\delta\in(0,\frac14]$, $\eta\in(0,1]$, a non-increasing $\bar\ell(\cdot)\to0$, a threshold $t_*$, and the \emph{predictable} good indicator $g_i=\1\{t_{i-1}\ge t_*\}\1\{\theta_i\in R(\delta)\}\1\{\bar\ell_i\le\bar\ell(t_{i-1})\}$. Inputs: (H3) $\P(\exists i:t_{i-1}\ge T,g_i=0)\le C_1e^{-c_1T}$ for $T\ge t_*$ (Theorem~\ref{thm:Cexp}, information concentration, Lemma~\ref{lem:E0}); (H4) $t_{i-1}\ge i\eps/(2\mu_{hi})$ except with probability $\le Ce^{-ci}$ (upper tail of $\Lambda_T$); (H5) uniformly bounded moments of the gap quantities. Then $\sum_i\P(g_i=0)^{1/2}\le C_b$, hence $\E\sum_i(1-g_i)\Phi_i\le C_B$ for every $\Phi_i\ge0$ with $\E\Phi_i^2\le C_\Phi$, independently of $n$.
\emph{(F1)} $R_i\le X_i:=v_i\T\Delta F_iv_i$ pathwise ($\psi(r)\le(r-1)^2/r$). \emph{(F2)} On $\{\theta_i\in R(\delta)\}$, $\psi(r)\le(\frac12+\delta)(r-1)^2/r$ so $\E[g_iR_i]\le(\frac12+\delta)\E[g_iX_i]$. \emph{(F3)} $\E[g_iH_i\mid\cG_{i-1}]\le g_i\frac{\bar\ell(t_{i-1})}{\mu_{lo}}\E[X_i\mid\cG_{i-1}]$ (Lemma~\ref{lem:E5d} and the compensator identity). \emph{(F4)} $\E[g_iQ_i\mid\cG_{i-1}]=g_iq_i$ and, with $m_i=2\ln(1/\bar\ell(t_{i-1}))$, $q_i\le(1+\delta)r(\bar x_i)\E[\log\det(F_i/F_{i-1})\mid\cG_{i-1}]+3\bar\ell^3\ln(e/\bar\ell)/\mu^*$ (Lemma~\ref{lem:logdet} with $\lambda_\theta\le(1+\delta)\lambda^*$), $\bar x_i=\bar\ell m_i/(c^-\mu^*)\to0$, so $r(\bar x_i)\le1+\delta$ for $t_{i-1}\ge t_*$.
Assembly: by Proposition~\ref{prop:reg}(c), $\E X_n\le P_0+\E G$ with $G=\sum g_i\T F_i^{-1}g_i\le\sum g_i\T F_{i-1}^{-1}g_i$. Bad gaps contribute $\le C_B/\lambda_{\min}(F_0)$. Good gaps: $g\T F^{-1}g\le(1+\eta)Q+(1+1/\eta)H$ and (F3)--(F4) give
\[
\E X_n\le P_0+(1+\eta)(1+\delta)^2\E\log\det\tfrac{F_n}{F_0}+C_\rho+\tfrac{C_B}{\lambda_{\min}(F_0)}+B\,\E X_n,\qquad B=(1+\tfrac1\eta)\tfrac{\bar\ell(t_*)}{\mu_{lo}} .
\]
Choose $t_*$ with $B\le\delta$. Then by (F1)--(F2), $\E R_T\le(\frac12+\delta)\E X_T+C_B$, and $\E\log\det(F_n/F_0)\le\log\det(F_0^{-1}(F_0+T\,I/c^-))+C=d\ln T+O(1)$ (concavity of $\log\det$, $\Delta F_i\preceq\Delta J_i/c^-$, deterministic horizon $T$). Hence $\E R_T\le(\frac12+\delta)(1-\delta)^{-1}(1+\eta)(1+\delta)^2\,d\ln T+C$. No exponential tail of $G$, restart or Ville inequality is used in the assembly; they enter only through (H3). Full details: \texttt{theoremR\_proof\_complete.md}.
\end{proof}
Inputs and limits: Theorem~\ref{thm:E3}/\ref{thm:Fconc} (rate constants $\sim10^5$ conservative), Assumption~\ref{ass:M} (certified by Proposition~\ref{prop:Mexact} for the whole experimental range of $\eps$; Corollary~\ref{cor:M} is the older Gamma-comparator certificate for $\eps\le10^{-4}$), the a-priori box (cost: Proposition~\ref{prop:boxcost}), $K=1$ here; general $K$: Theorem~\ref{thm:RK}. Observed slope of $\E R_n$ against $\ln T_n$: $\approx0.9$ vs $d/2=1$ (\Nt{}).

\subsection{Assumption M by Fourier quadrature (any $\eps$, any $k_0$)}\label{sec:Mexact}

\begin{proposition}[exact certificate for Assumption~\ref{ass:M}]\label{prop:Mexact}\Ps{}
Let $h$ be a trapezoid centred at $c$ with plateau half-width $p$, support half-width $s$, $\rho=s-p$, $A=(s+p)/2$. Its Fourier transform is $e^{itc}H(t)$ with $H(t)=\frac4\rho\frac{\sin(At)\sin(\rho t/2)}{t^2}$, $|H|\le4/(\rho t^2)$, and by Parseval
\[
\E_\pi h(X)=\frac1\pi\int_0^\infty H(t)\,\mathrm{Re}\big[e^{itc}\overline{\varphi(t)}\big]\dd t,\qquad X=\eps\lambda,
\]
with $\varphi$ the exact characteristic function of Section~\ref{sec:tv} ($\dd L/\dd x=\tfrac k2x/E_1(x)$). Since $|\varphi|$ is non-increasing (Lemma~\ref{lem:elem}(c)), the tail beyond $T$ is at most $|\varphi(T)|\,4/(\pi\rho T)$.
\end{proposition}
Evaluated by Gauss--Legendre panels (two different panel/node choices agree to $\le2\times10^{-9}$) with the explicit tail bound $\le5\times10^{-6}$ (Table~\ref{tab:Mexact}); validated against stationary simulation ($0.15,0.05$) and the Gamma-comparator limits ($0.164,0.0444$).
\begin{table}[h]\centering\small
\begin{tabular}{@{}lccccc@{}}\toprule
$k_0$ & $\eps=0.3$ & $0.1$ & $0.03$ & $0.01$ & $0.001$\\\midrule
4 & $.1527/.0432$ & $.1570/.0425$ & $.1587/.0423$ & $.1592/.0423$ & $.1594/.0422$\\
2 & $.1528/.0463$ & $.1602/.0450$ & $.1628/.0446$ & $.1636/.0445$ & $.1639/.0444$\\
1.2 & $.1363/.0418$ & $.1418/.0403$ & $.1445/.0399$ & $.1453/.0397$ & $.1456/.0397$\\
0.6 & $.0855/.0324$ & $.1056/.0306$ & $.1088/.0301$ & $.1094/.0300$ & --\\\bottomrule
\end{tabular}
\caption{$\E_\pi h_1/\E_\pi h_2$ on $I_1=0.5\,\mathrm{ES}$, $I_2=1.8\,\mathrm{ES}$.}\label{tab:Mexact}
\end{table}
Hence $m\ge0.032$ ($k_0=0.6$), $0.040$ ($1.2$), $0.042$ ($4$), $0.044$ ($2$) for all $\eps\le0.3$ tested; Assumption~\ref{ass:M} is certified at every $\eps$ used in the regret experiments, and no condition $k_0>1$ is needed. The earlier Gamma-comparator route (Corollary~\ref{cor:M}) is superseded for this purpose. Limits: floating-point quadrature with convergence check (not interval arithmetic), per parameter value, $K=1$ only.

\subsection{What the a-priori box costs}\label{sec:boxcost}

The box $\Theta_c\ni\theta^*$ enters only through $c^\pm$ (information level $\lambda_T=mT\lambda_Q/(4c^+L)$, constant $A$ in Prop.~\ref{prop:E5}), $\kappa_1=\max(1,1/c^--1)$ (so $x/\theta_M\propto x/c^-$ in Lemma~\ref{lem:E4}), $\mu_{lo}$ (leverage condition $\bar\ell\lesssim\mu_{lo}$, i.e.\ leverage burn-in $\propto c^+/\mu_{lo}$), and $P_0$. After the bootstrap $c^\pm\to1\pm\delta$, so the leading term is box-free.

\begin{proposition}[cost of the box]\label{prop:boxcost}\Ps{}
Theorem~\ref{thm:R} holds for any compact box containing $\theta^*$ with the same leading constant; the additive constant grows at most polynomially in the box, like $c^+\max(1,1/c^-)$ in $T_1$ and $c^+/\mu_{lo}$ in the leverage burn-in, times the dependence of $P_0$ on the prior weight $\rho_0$.
\end{proposition}
\Nt{} ($K=1$, $n=3\times10^4$, corner start $(s\mu^*,\alpha^*/s)$ vs truth start): extra regret $\propto s^{2.0\text{--}2.1}$ at $\rho_0=1$ (e.g.\ $\eps=0.3$: $4.6,80,2449$ at $s=5,20,100$), reduced by a weak prior (at $s=20$: $21.1,1.7,0.35$ for $\rho_0=1,0.1,0.01$); the cost is a prior-weight effect. In the regime $\eps=c/T$ the scale of $\mu^*$ cannot be learned (Theorem~\ref{thm:V}), so there the box is a genuine requirement; for fixed $\eps$ a two-stage scheme (wide box, then a $(1+o(1))$-box around the first-stage estimate) should remove it but is not proved (\Ot{}).

\subsection{General $K$}\label{sec:generalK}

Everything in Sections~\ref{sec:regret}.1--\ref{sec:regret}.5 except the two $K=1$ devices (Assumption~\ref{ass:M} and Lemma~\ref{lem:E0}) is dimension-free: Prop.~\ref{prop:reg} needs only convexity of the polytope $\Theta_c=\{\mu\in[\mu_{lo},\mu_{hi}],\alpha_k\in[\alpha_{lo,k},\alpha_{hi,k}],\sum\alpha_k/\beta_k\le1-\eps/2\}$ and non-expansiveness of the $F$-projection (a convex QP); Prop.~\ref{prop:ratio} holds with coordinate ratios; Lemmas~\ref{lem:score}--\ref{lem:logdet}, \ref{lem:E5a}--\ref{lem:E5d}, \ref{lem:E4}, \ref{lem:C1} and Prop.~\ref{prop:E5} use only $\lambda^*\ge\mu^*$ and the time-change structure, with $d=K+1$. Lemma~\ref{lem:E0} is replaced by Lemma~\ref{lem:expmoment}, while Lemma~\ref{lem:E1}(a) holds verbatim for every $K$ with $Z_c=\sum_k(\alpha_k/\beta_k)S_k$: $\eps\Lambda_T=\mu T+nM'_T-(Z_c(T)-Z_c(0))$ (integrated key identity of Proposition~\ref{prop:keyid}).

\begin{theorem}[concentration of the information without Assumption M]\label{thm:FK}\Ps{}
Let $\lambda_*=\sup_{\lambda_0}\lambda_{\min}(N(\lambda_0))>0$ (Theorem~\ref{thm:tangent}). Pathwise $J_T\succeq\int_0^Taa\T(2/\lambda_0-\lambda^*/\lambda_0^2)\dd t$, whose entries are time integrals of the cubic polynomials $f_{ij}(S)=a_ia_j(2/\lambda_0-\lambda(S)/\lambda_0^2)$ with $\E f_{ij}=N(\lambda_0)_{ij}$. Clipping the state at $s_{\max}$ (componentwise) makes $f_{ij}$ Lipschitz in $d_\ell$ with constant $\Lambda_f(s_{\max})=O(s_{\max}^2)$; on $\{\max_{t\le T}S\le s_{\max}\}$ (probability $\ge1-T^{-\kappa}$ for $s_{\max}=O(\eps^{-1}\ln T)$, Lemma~\ref{lem:expmoment}) the clipping is inactive, and $|\E\tilde f-\E f|$ is exponentially small. The Poisson-equation/Freedman argument ($\hat u$ is $\Lambda_f/\gamma$-Lipschitz by Lemma~\ref{lem:ellcontr}, $|\Delta\mathcal M|\le\Lambda_f/\gamma$, $\langle\mathcal M\rangle\le(\Lambda_f/\gamma)^2\Lambda_T$) and a union over the $(K+1)(K+2)/2$ entries with $\tau=\lambda_*/2$ give, with probability $\ge1-3(K+1)(K+2)e^{-c_fT}-T^{-\kappa}$, $c_f\asymp\tau^2\gamma^2\eps/(\Lambda_f(s_{\max})^2\mu)$:
\[
J_T\succeq\frac{T\lambda_*}{2}\,I_{K+1}.
\]
\end{theorem}
Thus Assumption~\ref{ass:M} (a statement about $\pi$-mass of $K+1$ well-placed regions, not certifiable for $K\ge2$ by the Fourier method) is replaced by the computable number $\lambda_L>0$. Leverage: $\bar\ell(t)\le\frac{4Lc^+}{t\lambda_L\mu^*}(1+Ks_{\max}^2)=O((\ln t)^2/(\eps^2t))$.

\begin{theorem}[expected regret, general $K$]\label{thm:RK}\Ps{}
Under a known polytope $\Theta_c\ni\theta^*$ (information level $\lambda_*/2$ of Theorem~\ref{thm:FK}): for every $\delta,\eta>0$, $\E R_n\le(\frac12+O(\delta+\eta))(K+1)\ln T_n+C_K$ with $C_K$ independent of $n$; $\limsup_n\E R_n/\ln T_n\le(K+1)/2$. The proof is that of Theorem~\ref{thm:R} with the replacements above.
\end{theorem}
With the tangent level the constants are far less pessimistic than with truncation: for $\beta=(1,2)$, $\mu=1.5$, $\eps=0.4$ (0.3) the information level is $\lambda_*/2\approx0.0126$ ($0.0132$) against $\approx10^{-6}$ for the truncation device, a gain of about $10^4$ in $T_1$; they remain far from sharp (the rate $c_f$ contains $\Lambda_f(s_{\max})^{-2}$). \Nt{} ($K=2$, $\beta=(1,2)$, $\mu^*=1.5$, projection by SLSQP active in $0.2\%$ of steps, 6 replicates, $n=4\times10^4$): $\E R_n=4.41,5.29,6.14,7.28,7.99$ ($\eps=0.4$) and $4.85,5.86,6.60,7.64,8.28$ ($\eps=0.3$) at $n/16,\dots,n$; slope in $\ln T$: mean $1.3$ (both), target $(K+1)/2=1.5$; slope in $\log\det F$: $0.43$ vs $0.5$. This supports the order, not the exact constant, at these horizons.

\subsection{A minimax lower bound for any algorithm}\label{sec:lb}

An algorithm defines a predictive intensity $\hat\lambda_i(u)\ge\hat\mu_i>0$ on each gap, hence a probability law $Q^{(n)}$ on records of $n$ events; the cumulative log-loss regret equals $\KL(P_\theta^{(n)}\|Q^{(n)})=\E R_n$.

\begin{proposition}[lower bound]\label{prop:LB}\Ps{}
Let $\Theta_0\subset\mathrm{int}\,\Theta_\eps$ be a fixed compact box ($d=K+1$) and suppose an estimator $\hat\theta_n$ satisfies (U): $|\Theta_0|^{-1}\!\int_{\Theta_0}\E_\theta|\hat\theta_n-\theta|^2\dd\theta\le C_1\ln n/n$. Then for every $Q^{(n)}$
\[
\sup_{\theta\in\Theta_0}\KL(P_\theta^{(n)}\|Q^{(n)})\ \ge\ \ln|\Theta_0|+\frac d2\ln\frac n{\ln n}-\frac d2\ln\frac{2\pi eC_1}d ,
\]
and $\frac d2\ln n-O(1)$ if (U) holds with $C_1/n$.
\end{proposition}
\begin{proof}
$\sup\ge$ average over $\Theta\sim\mathrm{Unif}(\Theta_0)$; $\E_\Theta\KL(P_\Theta\|Q)=I(\Theta;X)+\KL(\bar P\|Q)\ge I(\Theta;X)\ge I(\Theta;\hat\theta_n)=h(\Theta)-h(\Theta|\hat\theta_n)$, with $h(\Theta)=\ln|\Theta_0|$ and $h(\Theta|\hat\theta_n)\le h(\Theta-\hat\theta_n)\le\frac d2\ln(2\pi e\,\E|\Theta-\hat\theta_n|^2/d)$ (Gaussian maximum entropy, AM--GM).
\end{proof}
Hypothesis (U) is delivered, for fixed $\eps$, by Theorem~\ref{thm:Cexp} and the information level ($D_n\le C(d\ln n+x)$, $\lambda_{\min}(F_n)\ge cn$ w.h.p.), with constants continuous in the parameters, hence bounded on a compact $\Theta_0$ (stated, not proved). Together with Theorem~\ref{thm:R}: $\inf_Q\sup_{\Theta_0}\E R_n=\frac d2\ln n\,(1+o(1))$, so the leading constant $d/2$ is minimax optimal. The first-order term in time units is $\frac d2\ln T+\frac12\ln\det I$, and $\det I$ stays bounded as $\eps\to0$ for $k_0>1$ (Thm.~\ref{thm:A}), consistent with the absence of a $\ln(1/\eps)$ term in the experiments; an $\eps$-explicit bound for $\eps=\eps(n)\to0$ is open (\Ot{}).


\subsection{Matrix-covariance lower bound, its $\eps$-dependence and the two-point form}\label{sec:lbmatrix}

\begin{proposition}[two-point bound in Fisher form]\label{prop:twopoint}\Pt{}
Let $\theta,\theta'$ lie in a ratio set ($c^-\le\lambda_{\theta'}/\lambda_\theta\le c^+$) and let both laws start from the same observed state. Then, with $\tilde c^{+}=\max(1,c^+)$, $\tilde c^-=\min(1,c^-)$ and $u=\theta'-\theta$,
\[
\frac{T\,u\T\bar I_Tu}{2\tilde c^{+}}\ \le\ \KL(P_\theta^T\|P_{\theta'}^T)\ \le\ \frac{T\,u\T\bar I_Tu}{2(\tilde c^{-})^2},\qquad \bar I_T=\tfrac1T\E_\theta J_T\ \ (=I(\theta)\text{ for a stationary start}).
\]
\end{proposition}
\begin{proof}
$\KL(P_\theta^T\|P_{\theta'}^T)$ is the regret of the constant predictor $\theta'$ (Lemma~\ref{lem:regid}); apply Lemma~\ref{lem:regsand} and take expectations.
\end{proof}
The bound is direction-resolved ($u\T Iu\sim\eps/(\mu(k_0-1))$ along the ridge, $\sim1/\eps$ along the top eigenvector), unlike the kernel-agnostic $\delta^2T\mu(1-n)$ bound of the literature.

\begin{proposition}[matrix lower bound]\label{prop:LBmat}\Pt{}
Let $\Theta\sim\mathrm{Unif}(\Theta_0)$ and suppose a statistic $\hat\theta$ satisfies $\E_\theta[(\hat\theta-\theta)(\hat\theta-\theta)\T]\preceq(\kappa_T/T)A^{-1}$ on $\Theta_0$, $A\succ0$. Then for every predictor $Q$
\[
\sup_{\theta\in\Theta_0}\mathrm{Reg}_T\ \ge\ I(\Theta;X)\ \ge\ \ln|\Theta_0|+\tfrac d2\ln\frac{T}{2\pi e\kappa_T}+\tfrac12\ln\det A .
\]
\end{proposition}
\begin{proof}
As Proposition~\ref{prop:LB}, with the Gaussian maximum-entropy bound $h(\Theta-\hat\theta)\le\frac12\ln\big((2\pi e)^d\det\E[(\Theta-\hat\theta)(\Theta-\hat\theta)\T]\big)$ and monotonicity of $\det$ in the Loewner order.
\end{proof}

\begin{theorem}[lower bound with explicit $\eps$-dependence]\label{thm:LBeps}\Ps{}
For online Newton, $\hat\theta=\theta_{n+1}$ satisfies the hypothesis of Proposition~\ref{prop:LBmat} with $A=N(\theta)$ (the tangent minorant of $I$) and $\kappa_T=2c^+c_Dd\ln T$ on the good event of Theorem~\ref{thm:Cexp}. Hence for $\Theta_0=[\mu(1-w),\mu(1+w)]\times[\beta(1-\tfrac32\eps),\beta(1-\tfrac12\eps)]$ ($K=1$) and $T\eps\gg\ln T$
\[
\sup_{\Theta_0}\mathrm{Reg}_T\ \ge\ \ln(T\eps)+\ln(2w\mu\beta)+\tfrac12\ln\det N_*-\ln(2\pi e\kappa_T),
\]
and the upper bound of Theorem~\ref{thm:R} with a prior matched to $\Theta_0$ is $\ln(T\eps)+\ln(2w\mu\beta)+\tfrac12\ln\det I_*+O(1)$. For $k_0>1$ both $\det I_*$ and $\det N_*$ are bounded above and below uniformly in $\eps$, so the two bounds agree in their $\eps$-dependence up to $\tfrac12\ln\ln T+O(1)$; for $k_0<1$ the upper bound carries the additional $\tfrac{1-k_0}2\ln(1/\eps)$ of $\tfrac12\ln\det I$.
\end{theorem}


%==============================================================
\section{The stationary record: total-variation bounds and impossibility}\label{sec:tv}
%==============================================================

\begin{lemma}[TV from characteristic functions]\label{lem:LemmaF}\Pt{}
For densities $p,q$ with characteristic functions $\varphi_p,\varphi_q$ and $\Delta=\varphi_p-\varphi_q$: $\TV(p,q)\le\norm\Delta_{H^1}/(2\sqrt2)$, $\norm\Delta_{H^1}^2=\int(|\Delta|^2+|\Delta'|^2)\dd t$.
\end{lemma}
\begin{proof}
$\int|\rho|\le(\int(1+x^2)^{-1}\dd x)^{1/2}(\norm\rho_2^2+\norm{x\rho}_2^2)^{1/2}=\sqrt\pi\big((\norm\Delta_2^2+\norm{\Delta'}_2^2)/(2\pi)\big)^{1/2}$ by Plancherel.
\end{proof}

\paragraph{Setting ($K=1$, $\beta=1$, $k_0>1$).} $X=\eps\lambda$; $\varphi(t)=\E e^{itX}=\exp(L(t))$, $L(t)=-\beta\mu\int_0^{-i\eps t}s\,\dd s/c(s)$ (Lemma~\ref{lem:laplace} on the imaginary axis; $c$ has no zero in $\mathrm{Re}\,s\ge0$, $s\ne0$). By Lemma~\ref{lem:gammafact}'s algebra, $\varphi=G\tilde F$ with $G(t)=(1-iat)^{-k}$ (Gamma$(k,a)$, $a=\alpha^2/(2\beta)$) and $\tilde F=\exp(z)$, $z(x)=i\frac k2\int_0^x\varsigma(y)\dd y$, $x=\alpha\eps t$, where with $r=\eps/(1-\eps)$, $y=\alpha\tau$:
\[
\varsigma(y)=\alpha D_\eps(-iy/\alpha)=\frac{N(y)}{E_1(y)(r-iy/2)},\ \ N=(1-\cos y-\tfrac{y^2}2)+i(y-\sin y),\ \ E_1=-(1-\cos y)-i((1+r)y-\sin y).
\]
Note $k/(1+r)^2=k_0$ exactly (with $k=k_0/(1-\eps)^2$).

\begin{lemma}[elementary bounds]\label{lem:elem}\Pt{}
\begin{enumerate}[label=(\alph*),nosep]
\item $|\varsigma(y)|\le19$ for all $y\ge0$, $\eps\in(0,1)$. ($y\le1$: $|N|\le0.2084y^3$, $|E_1|\ge0.4583y^2$, $|r-iy/2|\ge y/2$, giving $\le0.91$. $y\ge1$: $\mathrm{Re}\,N\in[-y^2/2,0]$ because $1-\cos y\le y^2/2$, so $|N|\le y^2/2+y\le1.5y^2$; $|E_1|\ge y-\sin y\ge0.158y$; so $|\varsigma|\le18.99$.) True supremum: $1.1429$.
\item With $M_1:=19k\alpha/2$: $|\tilde F(t)|\le e^{M_1\eps|t|}$, $|\tilde F'|\le M_1\eps|\tilde F|$.
\item $|\varphi|$ is non-increasing on $t\ge0$: $\mathrm{Re}\,L'=-\frac k2x(1-\cos x)/|E_1|^2\le0$.
\item $|L'(t)|\le1.1k/t$ for $x\le1$ and $\le3.2k\alpha\eps$ for $x\ge1$ (from $|E_1|\ge0.4583x^2$, resp.\ $\ge0.158x$).
\item For $x\ge5$: $|E_1|^2\le(1+r)^2x^2(1+3/x)$.
\end{enumerate}
\end{lemma}

\begin{theorem}[Gamma comparator with explicit constants]\label{thm:Tprime}\Ps{}
Let $k_0>1$, $\eps\le\frac12$, $\Gamma_\theta=\mathrm{Gamma}(k,a/\eps)$ in $\lambda$-space. Then $\TV(\pi_\theta,\Gamma_\theta)\le\Phi(k_0,\eps)/(2\sqrt2)$, with $\Phi$ the explicit number built as follows. Put $t_1=1/(M_1\eps)$, $X_0=\max(5,6/(k_0-1))$, $t_*=X_0/(\alpha\eps)$, $q=k_0/(2(1+3/X_0))\ge k_0/(k_0+1)>\frac12$, $A=e(at_1)^{-k}$.
\begin{itemize}[nosep]
\item \emph{Head} $|t|\le t_1$: $|z|\le1$, $|\tilde F-1|\le(e-1)M_1\eps|t|$, $|\Delta|\le(e-1)M_1\eps|t||G|$, $|\Delta'|\le M_1\eps(k(e-1)+e)|G|$ (using $|G'|\le(k/|t|)|G|$). So head$^2\le(e-1)^2M_1^2\eps^2J+M_1^2\eps^2(k(e-1)+e)^2a^{-1}B(\tfrac12,k-\tfrac12)$, $J=\int_{|t|\le t_1}t^2(1+a^2t^2)^{-k}\dd t$ ($\le a^{-3}B(\frac32,k-\frac32)$ for $k>\frac32$).
\item \emph{Tail} $|t|>t_1$: by monotonicity $|\varphi|\le A$ on $[t_1,t_*]$; for $x\ge X_0$, (c),(e) and $\int_{X_0}^x(1-\cos y)/y\dd y\ge\ln(x/X_0)-2/X_0$ give $|\varphi(t)|\le Ae^{k_0/X_0}(t_*/t)^q$ for $t\ge t_*$. Integrating (two sides) gives $I_\varphi,I_{\varphi'}$ (using (d) for $\varphi'=\varphi L'$) and the Gamma tails $I_G=2a^{-2k}t_1^{1-2k}/(2k-1)$, $I_{G'}=2k^2a^{-2k}t_1^{-2k-1}/(2k+1)$.
\item $\Phi^2:=\mathrm{head}^2+2(I_\varphi+I_G)+2(I_{\varphi'}+I_{G'})$; $\Phi=O(\eps^{\min(1,k-1/2)})$.
\end{itemize}
\end{theorem}
\emph{Values} ($k_0=2$) of $\Phi/(2\sqrt2)$: $2.0\times10^{-2}$ ($\eps=10^{-4}$), $3.9\times10^{-3}$ ($3\times10^{-5}$), $1.0\times10^{-3}$ ($10^{-5}$), $8.6\times10^{-5}$ ($10^{-6}$), $8.4\times10^{-7}$ ($10^{-8}$); $\Phi$ exceeded the true $\norm{\varphi-G}_{H^1}$ in every test ($\eps$ from $0.1$ to $10^{-8}$) by factors $75$ to $7\times10^4$; vacuous at $\eps=0.1$. Other $k_0$ at $\eps=10^{-6}/10^{-8}/10^{-10}$: $k_0=1.2$: $0.10/4.0\times10^{-3}/1.6\times10^{-4}$; $k_0=3$: $1.3\times10^{-4}/1.3\times10^{-6}/1.3\times10^{-8}$.
The chain of inequalities is proved by hand; the numerical value of $\Phi$ is a floating-point quadrature, not interval arithmetic.

\begin{proposition}[two Gamma laws]\label{prop:twogamma}\Nt{}
For the ridge pair $\theta'=(\mu+\delta,\alpha(1-\eps\delta/\mu))$ the closed form $\KL(\Gamma(k,s)\|\Gamma(k',s'))=(k-k')\psi(k)-\ln\Gamma(k)+\ln\Gamma(k')+k'\ln(s'/s)+k(s-s')/s'$ was evaluated on a $\delta$-grid ($\mu=\beta=1$): $\KL/\delta^2\le0.30$ for $\delta\le0.3$, so $\TV(\Gamma_\theta,\Gamma_{\theta'})\le0.39\delta$ (Pinsker). A Hessian bound on the log-partition function would make this a proof; unwritten.
\end{proposition}

\begin{corollary}[stationary law on the ridge, $k_0=2$, $\mu=\beta=1$]\label{cor:TVridge}\Ps{}
For $\eps\le10^{-5}$ and $|\delta|\le0.3$: $\TV(\pi_\theta,\pi_{\theta'})\le2.4\times10^{-3}+0.39|\delta|$ (first term evaluated at $\delta=0,0.1,0.3$ from $\Phi(2,\eps)$ and $\Phi(2(1+\delta),\eps(1+(1-\eps)\delta))$).
\end{corollary}

\begin{corollary}[Assumption M certified]\label{cor:M}\Ps{}
For $k_0=2$, $\mu=\beta=1$: $\E_\pi h_j\ge\E_\Gamma h_j-\TV(\pi,\Gamma)$ with $\E_\Gamma h_1=0.164$, $\E_\Gamma h_2=0.0444$ (incomplete-gamma integrals); hence $m\ge0.0248$ for $\eps\le10^{-4}$ and $\ge0.0405$ for $\eps\le3\times10^{-5}$.
\end{corollary}

\begin{theorem}[impossibility for the stationary record, $K=1$, $k_0>1$]\label{thm:imp1}\Ps{}
With $\eps=c/T$, $\TV(Q_\theta,Q_{\theta'})\le\sqrt{\rho\,u\T I_Tu/2}+\TV(\pi_\theta,\pi_{\theta'})$ (data processing, Lemma~\ref{lem:KL}, Pinsker); by Theorem~\ref{thm:B} and Corollary~\ref{cor:TVridge} both terms are small for small $\delta$, so by Le Cam no estimator is consistent for $\mu$ from the stationary record with unobserved past. Explicit constants are available for the stationary-law term only in the range stated in Corollary~\ref{cor:TVridge}; the KL term keeps the constants of Theorem~\ref{thm:B}.
\end{theorem}

\paragraph{Any $K$ (route that avoids $\TV(\pi_\theta,\pi_{\theta'})$).}
Let $R_L$ be the law of the record on $[0,T]$ of the process started at time $-L$ from $S_{-L}\sim\pi_\theta$ and run with $\theta'$-dynamics, $L=K_0/\gamma$.
Step~1: $\KL(P_\theta|_{[-L,T]}\|R_L)\le(L+T)\rho\,u\T I(\theta)u\le\rho(1+\eps)^2C_E(K_0\tilde\nu+c)\delta^2$ (same initial law; Theorem~\ref{thm:E}); hence $\TV(Q_\theta,R_L)\le\delta[\rho(1+\eps)^2C_E(K_0\tilde\nu+c)/2]^{1/2}$.
Step~2 (superposition): from time $-L$ the intensity is $\mu'+g_s(t)+\sum_i\alpha'(t-t_i)$, $g_s(t)=\sum_k\alpha'_ks_ke^{-\beta_k(t+L)}$; by the Hawkes--Oakes cluster representation the process is the union of an empty-start process (shared) and independent ``ghost'' immigrants of intensity $g_s$ with their clusters. So $\TV(R_L,Q_{\theta'})\le\P(X_s\text{ hits }[0,T])+\P(X_{s'}\text{ hits }[0,T])$.
Step~3: ghost-hit probability via the survival function of a cluster.

\begin{lemma}[cluster survival, any $K$]\label{lem:S}\Pt{}
Let $q(r)=\P(\text{a cluster rooted at }0\text{ has an event after time }r)$. Then $q(s)=1-e^{-|Y(s)|}$, $|Y|=\sum_kY_k$, $Y'=-\mathrm{diag}(\beta)Y+\alpha(1-e^{-|Y|})$ (componentwise, with $\mathrm{diag}(\beta)$ acting on $Y$), $Y(0)=(\alpha_k/\beta_k)$; with $\zeta=\sum_kY_k/(\beta_k-\gamma)$: $\zeta'=-\gamma\zeta-r_e(|Y|)$, $r_e(x)=x-1+e^{-x}$. For $s\ge1/\gamma$,
$q(s)\le C'\gamma e^{-\gamma s}$, $C'=(B-\gamma)/(c_r(b-\gamma)^2(1-e^{-1}))$, $c_r=r_e(R_0)/R_0^2$, $R_0=(B-\gamma)/(b-\gamma)$.
\end{lemma}
\begin{proof}
Offspring of the root form a Poisson process of intensity $\sum\alpha_ke^{-\beta_kr}$; a sub-cluster rooted at $r\le s$ has no event after $s$ w.p.\ $1-q(s-r)$, any offspring after $s$ is itself an event. So $1-q(s)=\exp(-\bar\alpha(s)-\int_0^s\alpha(r)q(s-r)\dd r)$, whose exponent is $|Y|$. For $\zeta$ use $\sum\alpha_k/(\beta_k-\gamma)=1$ and $\beta_k/(\beta_k-\gamma)=1+\gamma/(\beta_k-\gamma)$. Since $(b-\gamma)\zeta\le|Y|\le(B-\gamma)\zeta<R_0$ and $r_e(x)/x^2$ is decreasing, $r_e(|Y|)\ge c\zeta^2$, so $v=1/\zeta$ satisfies $v'\ge\gamma v+c$.
\end{proof}
(The weaker first-moment bound loses a factor $\eps^{-2}$ and would force $L\approx(2/\gamma)\ln(1/\eps)$.) Integrating the ghost intensity against $q$ gives $\E\P(X_s\text{ hits})\le C'\mu e^{-\gamma L}(1+O(\eps\delta))+e^{-b(L-1/\gamma)}\mu/(b\eps)$.

\begin{theorem}[impossibility for the stationary record, any $K$, $k_K>1$]\label{thm:V}\Ps{}
For $\eps=c/T\le\eps_0$ and $0<\delta\le\delta_1$ (constants depend on the dictionary, $\mu$, $c$): $\TV(Q_\theta,Q_{\theta'})\le C\delta(\ln(1/\delta)+1)^{1/2}$ (take $K_0=\ln(1/\delta)$). Hence for $\delta$ small, by Le Cam, $\max_{\theta,\theta'}\P(|\hat\mu-\mu_i|\ge\delta/2)\ge\frac14$: $\mu$ is not consistently estimable from a stationary record with unobserved past, with or without $\alpha$ known.
\end{theorem}
Limits: needs $k_K>1$ and Theorem~\ref{thm:E} ($\eps_0$ not explicit); the cluster/superposition representation with a deterministic time-varying baseline is standard but its proof was not written; the $\ln(1/\delta)$ factor is likely an artefact; whether $\TV(\pi_\theta,\pi_{\theta'})=O(\delta)$ for $K\ge2$ is open. \Nt{}: flow vs cluster simulation (60\,000 clusters) agree to $0.5$--$1\%$; ghost-hit probability decays by $e^{-1}$ per unit of $K_0$ as predicted.
For the empty-start experiment the baseline information diverges for every $k$ (Proposition~\ref{prop:cold}), so the lower bound there is logarithmic, not inconsistency.
